%% Extraído de docs/TCC_documento/Textuais/Metodologia/00_Introducao.tex (linhas 106-178)
    \begin{tikzpicture}[
        font=\small,
        flowbox/.style={draw, rounded corners, fill=white, align=center, minimum width=4.8cm, minimum height=0.9cm, text width=4.6cm, font=\footnotesize},
        small/.style={draw, rounded corners, fill=white, align=center, minimum width=2cm, minimum height=0.6cm, text width=1.9cm, font=\scriptsize},
        coltitle/.style={align=center, font=\bfseries},
        arrow/.style={-{Latex}, thick},
        qmark/.style={draw, diamond, fill=black!85, text=white, font=\bfseries\tiny, align=center, minimum width=0.9cm, minimum height=0.6cm, inner sep=0pt}
    ]
        % Faixas de fase (fundo), horizontais, cobrindo as três colunas
        \fill[green!45!black!25]  (-8.9,-0.7) rectangle (8.9,5.6);
        \fill[orange!55!black!20] (-8.9,-3.4) rectangle (8.9,-0.7);
        \fill[red!55!black!20]    (-8.9,-7)   rectangle (8.9,-3.4);

        \node[qmark] at (-9.6,5.6)  {Q1};
        \node[qmark] at (-9.6,-0.7) {Q2};
        \node[qmark] at (-9.6,-3.4) {Q3};

        % --- Coluna esquerda: Hardware (PCB) ---
        \node[coltitle] at (-6.0,6) {\textit{Hardware} (PCB)};
        \node[flowbox] (consid) at (-6.0,5.0) {Segurança, alimentação,\\ETH e micro};
        \node[flowbox] (cond)   at (-6.0,3.5) {Condicionamento\\das fases};
        \node[flowbox] (readeq) at (-6.0,2.0) {Readequação\\dos componentes};
        \node[flowbox] (restr)  at (-6.0,0.5) {Restrições físicas};
        \node[small]   (esqpcb) at (-7.15,-1.6) {Esquemático};
        \node[small]   (proto)  at (-4.85,-1.6) {Protótipo};
        \node[flowbox] (valpcb) at (-6.0,-4.5) {Validação das placas};

        \draw[arrow] (consid) -- (cond);
        \draw[arrow] (cond) -- (readeq);
        \draw[arrow] (readeq) -- (restr);
        \draw[arrow] (restr) -- (esqpcb);
        \draw[arrow] (restr) -- (proto);
        \draw[arrow] (esqpcb) -- (valpcb);
        \draw[arrow] (proto) -- (valpcb);
        \draw[-{Latex}, thick, dashed] (valpcb.west) to[bend right=35] node[right, pos=0.5, font=\scriptsize, align=center]{se reprovado,\\refatora} (esqpcb.south west);

        % --- Coluna central: Firmware (Código) ---
        \node[coltitle] at (0,6) {\textit{Firmware} (Código)};
        \node[flowbox] (cod)    at (0,5.1)  {Código};
        \node[flowbox] (dev)    at (0,3.9)  {Compra \textit{dev boards}};
        \node[flowbox] (zephyr) at (0,2.7)  {Criação do projeto em Zephyr};
        \node[flowbox] (amb)    at (0,1.5)  {Criação de ambiente de teste};
        \node[flowbox] (vcore)  at (0,-1.5) {Validação Vcore};
        \node[flowbox] (valsys) at (0,-4.5) {Validação System\\e Application};

        \draw[arrow] (cod) -- (dev);
        \draw[arrow] (dev) -- (zephyr);
        \draw[arrow] (zephyr) -- (amb);
        \draw[arrow] (amb) -- (vcore);
        \draw[arrow] (vcore) -- (valsys);

        % --- Coluna direita: Motor e observador ---
        \node[coltitle] at (6.0,6) {Motor e observador};
        \node[flowbox] (bldc) at (6.0,5.1) {Motor BLDC};
        \node[flowbox] (par)  at (6.0,3.9) {Parâmetros};
        \node[flowbox] (sim)  at (6.0,2.7) {Simulações};
        \node[flowbox] (obs)  at (6.0,1.5) {Observador};
        \node[flowbox] (dbt)  at (6.0,-1.5) {Teste simples\\com \textit{dev board}};
        \node[flowbox] (emb)  at (6.0,-4.5) {Observador embarcado\\na placa};

        \draw[arrow] (bldc) -- (par);
        \draw[arrow] (par) -- (sim);
        \draw[arrow] ([xshift=-6mm]sim.south) -- ([xshift=-6mm]obs.north);
        \draw[arrow] ([xshift=6mm]obs.north) -- ([xshift=6mm]sim.south);
        \draw[arrow] (obs) -- (dbt);
        \draw[arrow] (dbt) -- (emb);

        % Convergência
        \node[flowbox, fill=black!85, text=white] (placa_ok) at (0,-6.3) {Placa de medição completa};
        \draw[arrow] (valpcb.south) |- (placa_ok.west);
        \draw[arrow] (valsys) -- (placa_ok);
        \draw[arrow] (emb.south) |- (placa_ok.east);
    \end{tikzpicture}%
